\section{What a closure problem is}
\label{sec:problem}

\subsection{The abstract statement}

Suppose a system has a full state $X$ obeying $\dot X = F(X)$, and we only
want to track a reduced state $x = \cP X$: an average, a few moments, a
coarse-grained field. Applying $\cP$ to the dynamics gives
\begin{equation}
  \dot x = \cP F(X).
  \label{eq:unclosed}
\end{equation}
This equation is exact. Its right-hand side depends on the full state $X$,
not only on $x$, and that is the closure problem. A \emph{closure} is a rule
that replaces $\cP F(X)$ by something computable from $x$ alone, possibly
from its gradients or its history:
\begin{equation}
  \dot x = f(x) + \underbrace{\bigl[\cP F(X) - f(x)\bigr]}_{\text{closure defect } d} .
  \label{eq:defect}
\end{equation}
The bracket is what the closure gets wrong. It is a function of the full
state, so it can be measured only where the full state is available: on
reference trajectories or reference ensembles.

Two different objects are called ``the reduced equation'', and it helps to
keep them apart. Equation~\eqref{eq:unclosed} is an identity obtained by
averaging, filtering, taking moments or projecting. The closed model
$\dot x = f(x)$ is an approximation. The modelling enters only in the step
from the first to the second.

\subsection{Two ways to measure the error}

Let $X(t)$ be a reference solution, and let $x_c(t)$ solve the closed model
from the same reduced initial data.
\begin{itemize}
  \item The \textbf{a priori} error is the defect on the reference solution,
    $d(t) = \cP F(X(t)) - f(\cP X(t))$.
  \item The \textbf{a posteriori} (embedded) error is
    $e(t) = x_c(t) - \cP X(t)$.
\end{itemize}
Subtracting the two evolution equations gives
$\dot e = f(x_c) - f(\cP X) - d(t)$. If $f$ is Lipschitz with constant $L$
then by Gr\"onwall's inequality
\begin{equation}
  |e(t)| \le \int_0^t e^{L(t-s)}\,|d(s)|\,\dd s .
  \label{eq:gronwall}
\end{equation}
This inequality says three things. A small defect controls the embedded
error for times of order $1/L$ and no longer, unless the closed model is
contractive. The defect that matters is the one along the reference
trajectory, while a closed model run for a long time visits states the
reference never produced. And the bound is one-directional: a closure can
have a large pointwise defect and still give good embedded statistics, as
eddy-viscosity models do (Section~\ref{sec:turb}). These are the reasons the
benchmark requires an embedded-solver test next to every a priori test.

\subsection{Five items that describe any scale transition}

Every transition in this project is described by five items. Later sections
state them for each closure.
\begin{center}
\renewcommand{\arraystretch}{1.25}
\begin{tabularx}{\textwidth}{@{}P{3.4cm}Y@{}}
\toprule
Exact hierarchy & The chain of exact equations in which each retained
  quantity is coupled to one that is not retained. \\
Limiting process & The parameter that is sent to zero or infinity, and the
  scaling of time and space that goes with it. \\
Retained state & What the reduced model evolves. \\
Discarded information & What $\cP$ removes: velocity-space structure,
  correlations, fluctuations, phases, fast fields, history. \\
Closure defect & The bracket in \eqref{eq:defect}, with a statement of how
  it is measured. \\
\bottomrule
\end{tabularx}
\end{center}

\subsection{Where the problem comes from}

The problem has two causes, and usually both are present.
\begin{itemize}
  \item \textbf{Nonlinearity.} The average of a product is not the product of
    averages. Averaging a quadratic term produces a second moment, whose
    equation contains a third.
  \item \textbf{Transport.} Streaming or advection couples each retained
    quantity to one that is not retained. Multiplying
    $\bv\cdot\nabla f$ by $\bv^k$ and integrating always produces a moment
    of order $k+1$, even when the equation is linear.
\end{itemize}
A third situation is the elimination of whole degrees of freedom: fast
variables, or fields carried by another species. It is the same problem with
a different projection, and Section~\ref{sec:memory} treats it in general.

\subsection{A worked example with every feature}
\label{sec:fastslow}

\begin{example}[A slow variable driven by a fast one]
\label{ex:fastslow}
Let
\begin{equation}
  \dot x = a\,x + b\,y, \qquad
  \varepsilon\,\dot y = c\,x - y ,
  \label{eq:fastslow}
\end{equation}
with $\varepsilon>0$ small. Retain $x$ and discard $y$. The second equation
can be solved exactly,
\[
  y(t) = e^{-t/\varepsilon}\,y(0)
       + \frac{c}{\varepsilon}\int_0^t e^{-s/\varepsilon}\,x(t-s)\,\dd s ,
\]
so the exact reduced equation is
\begin{equation}
  \dot x(t) = a\,x(t)
   + \frac{bc}{\varepsilon}\int_0^t e^{-s/\varepsilon}\,x(t-s)\,\dd s
   + b\,e^{-t/\varepsilon}\,y(0) .
  \label{eq:fastslow-exact}
\end{equation}
The three terms are an instantaneous term, a \emph{memory} term, and a term
that depends on the discarded initial data. Nothing has been approximated.

Three closures follow from expanding the memory integral for slowly varying
$x$, using $x(t-s) = x(t) - s\,\dot x(t) + \dots$:
\begin{align}
  \text{zeroth order (Markov):}\quad & \dot x = (a+bc)\,x, \nonumber\\
  \text{first order:}\quad & \dot x = (a+bc)(1-\varepsilon bc)\,x, \\
  \text{first order, rational:}\quad
     & (1+\varepsilon bc)\,\dot x = (a+bc)\,x . \nonumber
\end{align}
The second line substitutes the zeroth-order equation for $\dot x$ in the
correction. The third keeps $\dot x$ as an unknown. They agree to
$O(\varepsilon)$ and differ at $O(\varepsilon^2)$.

For $a=0$, $b=1$, $c=-1$ the exact slow eigenvalue is
$\lambda = (-1+\sqrt{1-4\varepsilon})/(2\varepsilon)
= -1-\varepsilon-2\varepsilon^2-\dots$:
\begin{center}
\begin{tabular}{rrrrr}
\toprule
$\varepsilon$ & exact & Markov & first order & rational \\
\midrule
0.2   & $-1.3820$ & $-1$ & $-1.2000$ & $-1.2500$ \\
0.1   & $-1.1270$ & $-1$ & $-1.1000$ & $-1.1111$ \\
0.05  & $-1.0557$ & $-1$ & $-1.0500$ & $-1.0526$ \\
0.025 & $-1.0263$ & $-1$ & $-1.0250$ & $-1.0256$ \\
\bottomrule
\end{tabular}
\end{center}
The Markov error is $O(\varepsilon)$ and the first-order errors are
$O(\varepsilon^2)$.
\end{example}

The example contains the whole subject in miniature.
\begin{itemize}
  \item The exact reduced dynamics has memory. A local closure is obtained
    only by assuming that the memory kernel decays faster than $x$ changes.
    That is a statement about a small parameter, here $\varepsilon$, in
    kinetic theory the Knudsen number.
  \item The closure is \emph{asymptotic}: the kernel
    $\varepsilon^{-1}e^{-s/\varepsilon}$ tends to a delta function, and the
    corrections are a series in $\varepsilon$.
  \item The discarded initial data matter for a time of order $\varepsilon$
    (an initial layer). If they are unknown they act as noise.
  \item The same order of accuracy can be reached by formulas that behave
    very differently when $\varepsilon$ is not small. The polynomial form
    changes sign at $\varepsilon bc = 1$. The rational form does not. The
    same distinction separates the Burnett equations from their
    regularizations (Section~\ref{sec:burnett}).
  \item Keeping $y$ as a variable removes the closure problem at the cost of
    a larger, stiffer system. Extended-moment models make this trade
    (Section~\ref{sec:moment}).
\end{itemize}

\subsection{The kinds of closure}
\label{sec:kinds}

Closures differ in what the unclosed quantity is allowed to depend on.
\begin{center}
\renewcommand{\arraystretch}{1.25}
\begin{tabularx}{\textwidth}{@{}P{3.3cm}P{5.3cm}Y@{}}
\toprule
Kind & Form & Examples \\
\midrule
Algebraic & $u = g(x)$ at the same point and time
  & equation of state; Euler ($\boldsymbol\pi=\bq=0$); double-adiabatic
    laws; Eddington factor \\
Gradient & $u = -D(x)\,\nabla x$
  & Navier--Stokes--Fourier; eddy viscosity; Braginskii; diffusion
    approximation \\
Nonlocal in space & $u(\bx) = \int W(\bx,\bx')\,x(\bx')\,\dd\bx'$
  & Hammett--Perkins; nonlocal electron heat flux \\
Memory & $u(t) = \int_0^t K(s)\,x(t-s)\,\dd s$
  & Mori--Zwanzig; viscoelastic models \\
Extended state & evolve $u$ with its own closed equation
  & Grad; Reynolds-stress transport; $k$--$\varepsilon$; Cattaneo \\
Stochastic & $u = g(x) + \text{noise}$
  & Langevin models; stochastic subgrid forcing \\
\bottomrule
\end{tabularx}
\end{center}
These are not exclusive. An extended-state model is a memory closure with an
exponential kernel (Section~\ref{sec:embedding}), and for linear problems
every row is a choice of how to approximate one function of wavenumber and
frequency (Section~\ref{sec:between}).
